\documentclass[10pt,a4paper]{amsart}
\usepackage{amsmath,amssymb,amsthm,graphicx}
% Retained minimal preamble for the Stacks Project chapter flat.tex (More on Flatness).
% Verbatim excerpt of the chapter preamble: Xy-pic setup (native preamble lines 14-18) and the native
% theorem declarations with the mathematical macros (native preamble lines 154-217). Copied mechanically;
% nothing added, removed or redefined. The wrapper class and the amsmath/amssymb/amsthm/graphicx base
% packages are supplied by the assembler, so the class/externaldocument part of the native preamble is omitted.
% For commutative diagrams we use Xy-pic
\usepackage[all]{xy}

% We use 2cell for 2-commutative diagrams.
\xyoption{2cell}
% Theorem environments.
%
\theoremstyle{plain}
\newtheorem{theorem}[subsection]{Theorem}
\newtheorem{proposition}[subsection]{Proposition}
\newtheorem{lemma}[subsection]{Lemma}

\theoremstyle{definition}
\newtheorem{definition}[subsection]{Definition}
\newtheorem{example}[subsection]{Example}
\newtheorem{exercise}[subsection]{Exercise}
\newtheorem{situation}[subsection]{Situation}

\theoremstyle{remark}
\newtheorem{remark}[subsection]{Remark}
\newtheorem{remarks}[subsection]{Remarks}

\numberwithin{equation}{subsection}

% Macros
%
\def\lim{\mathop{\mathrm{lim}}\nolimits}
\def\colim{\mathop{\mathrm{colim}}\nolimits}
\def\Spec{\mathop{\mathrm{Spec}}}
\def\Hom{\mathop{\mathrm{Hom}}\nolimits}
\def\Ext{\mathop{\mathrm{Ext}}\nolimits}
\def\SheafHom{\mathop{\mathcal{H}\!\mathit{om}}\nolimits}
\def\SheafExt{\mathop{\mathcal{E}\!\mathit{xt}}\nolimits}
\def\Sch{\mathit{Sch}}
\def\Mor{\mathop{\mathrm{Mor}}\nolimits}
\def\Ob{\mathop{\mathrm{Ob}}\nolimits}
\def\Sh{\mathop{\mathit{Sh}}\nolimits}
\def\NL{\mathop{N\!L}\nolimits}
\def\CH{\mathop{\mathrm{CH}}\nolimits}
\def\proetale{{pro\text{-}\acute{e}tale}}
\def\etale{{\acute{e}tale}}
\def\QCoh{\mathit{QCoh}}
\def\Ker{\mathop{\mathrm{Ker}}}
\def\Im{\mathop{\mathrm{Im}}}
\def\Coker{\mathop{\mathrm{Coker}}}
\def\Coim{\mathop{\mathrm{Coim}}}

% Boxtimes
%
\DeclareMathSymbol{\boxtimes}{\mathbin}{AMSa}{"02}

%
% Macros for moduli stacks/spaces
%
\def\QCohstack{\mathcal{QC}\!\mathit{oh}}
\def\Cohstack{\mathcal{C}\!\mathit{oh}}
\def\Spacesstack{\mathcal{S}\!\mathit{paces}}
\def\Quotfunctor{\mathrm{Quot}}
\def\Hilbfunctor{\mathrm{Hilb}}
\def\Curvesstack{\mathcal{C}\!\mathit{urves}}
\def\Polarizedstack{\mathcal{P}\!\mathit{olarized}}
\def\Complexesstack{\mathcal{C}\!\mathit{omplexes}}
% \Pic is the operator that assigns to X its picard group, usage \Pic(X)
% \Picardstack_{X/B} denotes the Picard stack of X over B
% \Picardfunctor_{X/B} denotes the Picard functor of X over B
\def\Pic{\mathop{\mathrm{Pic}}\nolimits}
\def\Picardstack{\mathcal{P}\!\mathit{ic}}
\def\Picardfunctor{\mathrm{Pic}}
\def\Deformationcategory{\mathcal{D}\!\mathit{ef}}

% --- appended 2026-10-12 (Chief): the project own definitions of the reference/slogan/history
% display environments, copied mechanically from the shared Stacks preamble lines 8-12
% (\usepackage{verbatim} plus the three \newenvironment lines, which the chapter preamble keeps
% verbatim from the project preamble). Nothing added, removed or redefined: these environments
% carry no mathematical content and the project itself typesets them as comments. Without these
% three lines a counted statement whose own text uses them fails with Environment slogan undefined.
\usepackage{verbatim}
\newenvironment{reference}{\comment}{\endcomment}
\newenvironment{slogan}{\comment}{\endcomment}
\newenvironment{history}{\comment}{\endcomment}
\usepackage{cleveref}
\crefname{originaldefinition}{original definition}{original definitions}
\Crefname{originaldefinition}{Original definition}{Original definitions}
\makeatletter
\newlabel{algebra-definition-henselization}{{04GQ}{source}{}{sourcepointer-1}{}}
\newlabel{algebra-definition-henselization@cref}{{[originaldefinition][0][]04GQ}{[1][0][]source}{}{}{}}
\makeatother
\crefname{originallemma}{original lemma}{original lemmas}
\Crefname{originallemma}{Original lemma}{Original lemmas}
\makeatletter
\newlabel{derived-lemma-subcategory-left-resolution}{{05T7}{source}{}{sourcepointer-2}{}}
\newlabel{derived-lemma-subcategory-left-resolution@cref}{{[originallemma][0][]05T7}{[1][0][]source}{}{}{}}
\makeatother
\crefname{originallemma}{original lemma}{original lemmas}
\Crefname{originallemma}{Original lemma}{Original lemmas}
\makeatletter
\newlabel{lemma-flat-up-down-henselization}{{05VL}{source}{}{sourcepointer-3}{}}
\newlabel{lemma-flat-up-down-henselization@cref}{{[originallemma][0][]05VL}{[1][0][]source}{}{}{}}
\makeatother
\crefname{originallemma}{original lemma}{original lemmas}
\Crefname{originallemma}{Original lemma}{Original lemmas}
\makeatletter
\newlabel{more-algebra-lemma-dumb-properties-henselization}{{07QM}{source}{}{sourcepointer-4}{}}
\newlabel{more-algebra-lemma-dumb-properties-henselization@cref}{{[originallemma][0][]07QM}{[1][0][]source}{}{}{}}
\makeatother
\crefname{originallemma}{original lemma}{original lemmas}
\Crefname{originallemma}{Original lemma}{Original lemmas}
\makeatletter
\newlabel{more-algebra-lemma-flat-descent-tor-amplitude}{{068S}{source}{}{sourcepointer-5}{}}
\newlabel{more-algebra-lemma-flat-descent-tor-amplitude@cref}{{[originallemma][0][]068S}{[1][0][]source}{}{}{}}
\makeatother
\crefname{originallemma}{original lemma}{original lemmas}
\Crefname{originallemma}{Original lemma}{Original lemmas}
\makeatletter
\newlabel{more-algebra-lemma-flat-push-tor-amplitude}{{066J}{source}{}{sourcepointer-6}{}}
\newlabel{more-algebra-lemma-flat-push-tor-amplitude@cref}{{[originallemma][0][]066J}{[1][0][]source}{}{}{}}
\makeatother
\crefname{originallemma}{original lemma}{original lemmas}
\Crefname{originallemma}{Original lemma}{Original lemmas}
\makeatletter
\newlabel{more-algebra-lemma-last-one-flat}{{0653}{source}{}{sourcepointer-7}{}}
\newlabel{more-algebra-lemma-last-one-flat@cref}{{[originallemma][0][]0653}{[1][0][]source}{}{}{}}
\makeatother
\crefname{originallemma}{original lemma}{original lemmas}
\Crefname{originallemma}{Original lemma}{Original lemmas}
\makeatletter
\newlabel{more-algebra-lemma-tor-amplitude}{{0654}{source}{}{sourcepointer-8}{}}
\newlabel{more-algebra-lemma-tor-amplitude@cref}{{[originallemma][0][]0654}{[1][0][]source}{}{}{}}
\makeatother
\title{Tor amplitude is the same over a local ring, over its henselization and over its strict henselization}
\author{The Stacks Project Authors}
\date{stacks-project:flat.tex:local-snapshot-20261011; GNU FDL 1.2 or later}
\begin{document}
\maketitle
\noindent\textit{Editorial note.} All delivered mathematical text is the author's own, copied from the pinned Stacks Project chapter file flat.tex; no statement or proof step was rewritten, reordered or removed.\par\medskip
\noindent\textit{Editorial note.} Preamble provenance: the chapter's own retained minimal preamble (sources/stacks-flat/retained-preamble.tex, lines 1-85) is a verbatim excerpt of the native Xy-pic setup and of the native theorem declarations and mathematical macros, followed by the project's own reference/slogan/history display-environment definitions copied mechanically from the shared preamble; the wrapper class and the amsmath/amssymb/amsthm/graphicx base packages are supplied by the assembler, and extra\_packages is left empty because everything needed is already in the retained preamble. No macro was invented or redefined.\par\medskip
\noindent\textit{Editorial note.} Same-chapter pointers are located at their \textbackslash{}label lines inside the pinned chapter file; cross-chapter pointers are located in the project's own tag list staged in the snapshot (sources/stacks-chow/tags.txt), at the verbatim lines carrying their canonical tags.\par\medskip
\noindent\textit{Allowed standard prerequisites (editorial).}\begin{itemize}
\item Henselization and strict henselization, and faithful flatness of the local ring maps. Algebra, label algebra-definition-henselization (statement, proof line 248); More on Algebra, label more-algebra-lemma-dumb-properties-henselization (proof line 274). Fixes the rings considered and gives that the local ring maps from the local ring to its henselization and strict henselization are faithfully flat, which yields the easy implications by pushforward of tor amplitude.
\item Pushforward and descent of tor amplitude along faithfully flat ring maps. More on Algebra, labels more-algebra-lemma-flat-push-tor-amplitude (proof line 280) and more-algebra-lemma-flat-descent-tor-amplitude (proof line 283). The first gives the implications from the strict henselization down to the henselization and to the local ring, and the second gives the equivalences between the cases (6), (5) and (3) by descent.
\item Left resolutions by free modules, flatness of a last cokernel and the flat up and down lemma. Derived Categories, label derived-lemma-subcategory-left-resolution (proof line 293); More on Algebra, labels more-algebra-lemma-last-one-flat (proof line 296) and more-algebra-lemma-tor-amplitude (proof line 304); Same chapter flat.tex, label lemma-flat-up-down-henselization (proof line 298). Produces the bounded above complex of free modules representing the object, identifies a cokernel of its differentials as a flat module over the local ring, transports flatness to the strict henselization, and thereby exhibits a complex of flat modules representing the object, which is tor amplitude.
\end{itemize}\medskip
\ifdefined\hypertarget\hypertarget{sourcepointer-1}{}\fi
\noindent\textit{Source reference.} Original definition 04GQ: sources/stacks-chow/tags.txt, source lines 5841--5841. Navigation pointer for \textbackslash{}ref\{algebra-definition-henselization\} used inside the counted span (line 248); the argument that consumes it is retained in the span.\par\medskip
\ifdefined\hypertarget\hypertarget{sourcepointer-2}{}\fi
\noindent\textit{Source reference.} Original lemma 05T7: sources/stacks-chow/tags.txt, source lines 7471--7471. Navigation pointer for \textbackslash{}ref\{derived-lemma-subcategory-left-resolution\} used inside the counted span (line 293); the argument that consumes it is retained in the span.\par\medskip
\ifdefined\hypertarget\hypertarget{sourcepointer-3}{}\fi
\noindent\textit{Source reference.} Original lemma 05VL: sources/stacks-flat/pinned-original-source.tex, source lines 183--183. Navigation pointer for \textbackslash{}ref\{lemma-flat-up-down-henselization\} used inside the counted span (line 298); the argument that consumes it is retained in the span.\par\medskip
\ifdefined\hypertarget\hypertarget{sourcepointer-4}{}\fi
\noindent\textit{Source reference.} Original lemma 07QM: sources/stacks-chow/tags.txt, source lines 9837--9837. Navigation pointer for \textbackslash{}ref\{more-algebra-lemma-dumb-properties-henselization\} used inside the counted span (line 274); the argument that consumes it is retained in the span.\par\medskip
\ifdefined\hypertarget\hypertarget{sourcepointer-5}{}\fi
\noindent\textit{Source reference.} Original lemma 068S: sources/stacks-chow/tags.txt, source lines 8016--8016. Navigation pointer for \textbackslash{}ref\{more-algebra-lemma-flat-descent-tor-amplitude\} used inside the counted span (line 283); the argument that consumes it is retained in the span.\par\medskip
\ifdefined\hypertarget\hypertarget{sourcepointer-6}{}\fi
\noindent\textit{Source reference.} Original lemma 066J: sources/stacks-chow/tags.txt, source lines 7938--7938. Navigation pointer for \textbackslash{}ref\{more-algebra-lemma-flat-push-tor-amplitude\} used inside the counted span (line 280); the argument that consumes it is retained in the span.\par\medskip
\ifdefined\hypertarget\hypertarget{sourcepointer-7}{}\fi
\noindent\textit{Source reference.} Original lemma 0653: sources/stacks-chow/tags.txt, source lines 7887--7887. Navigation pointer for \textbackslash{}ref\{more-algebra-lemma-last-one-flat\} used inside the counted span (line 296); the argument that consumes it is retained in the span.\par\medskip
\ifdefined\hypertarget\hypertarget{sourcepointer-8}{}\fi
\noindent\textit{Source reference.} Original lemma 0654: sources/stacks-chow/tags.txt, source lines 7888--7888. Navigation pointer for \textbackslash{}ref\{more-algebra-lemma-tor-amplitude\} used inside the counted span (line 304); the argument that consumes it is retained in the span.\par\medskip
\begin{lemma}
\label{lemma-tor-amplitude-up-down-henselization}
Let $S$ be a scheme and $s \in S$ a point. Denote $\mathcal{O}_{S, s}^h$
(resp.\ $\mathcal{O}_{S, s}^{sh}$) the henselization (resp.\ strict
henselization), see
Algebra, Definition \ref{algebra-definition-henselization}.
Let $M^{sh}$ be an object of $D(\mathcal{O}_{S, s}^{sh})$.
Let $a, b \in \mathbf{Z}$.
The following are equivalent
\begin{enumerate}
\item $M^{sh}$ has tor amplitude in $[a, b]$ over $\mathcal{O}_{S, s}$,
\item $M^{sh}$ has tor amplitude in $[a, b]$ over $\mathcal{O}_{S, s}^h$, and
\item $M^{sh}$ has tor amplitude in $[a, b]$ over $\mathcal{O}_{S, s}^{sh}$.
\end{enumerate}
If $M^{sh} =
M^h \otimes_{\mathcal{O}_{S, s}^h}^\mathbf{L} \mathcal{O}_{S, s}^{sh}$
for $M^h \in D(\mathcal{O}_{S, s}^h)$ this is also equivalent to
\begin{enumerate}
\item[(4)] $M^h$ has tor amplitude in $[a, b]$ over $\mathcal{O}_{S, s}$, and
\item[(5)] $M^h$ has tor amplitude in $[a, b]$ over $\mathcal{O}_{S, s}^h$.
\end{enumerate}
If $M^h = M \otimes_{\mathcal{O}_{S, s}}^\mathbf{L} \mathcal{O}_{S, s}^h$
for $M \in D(\mathcal{O}_{S, s})$
this is also equivalent to
\begin{enumerate}
\item[(6)] $M$ has tor amplitude in $[a, b]$ over $\mathcal{O}_{S, s}$.
\end{enumerate}
\end{lemma}

\begin{proof}
By More on Algebra, Lemma
\ref{more-algebra-lemma-dumb-properties-henselization}
the local ring maps
$\mathcal{O}_{S, s} \to \mathcal{O}_{S, s}^h \to \mathcal{O}_{S, s}^{sh}$
are faithfully flat.
Hence (3) $\Rightarrow$ (2) $\Rightarrow$ (1) and
(5) $\Rightarrow$ (4) follow from
More on Algebra, Lemma \ref{more-algebra-lemma-flat-push-tor-amplitude}.
By faithful flatness the equivalences (6) $\Leftrightarrow$ (5) and
(5) $\Leftrightarrow$ (3) follow from
More on Algebra, Lemma \ref{more-algebra-lemma-flat-descent-tor-amplitude}.
Thus it suffices to show that
(1) $\Rightarrow$ (3), (2) $\Rightarrow$ (3), and (4) $\Rightarrow$ (5).

\medskip\noindent
Assume (1). In particular $M^{sh}$ has vanishing cohomology
in degrees $< a$ and $> b$. Hence we can represent $M^{sh}$ by a complex
$P^\bullet$ of free $\mathcal{O}_{X, x}^{sh}$-modules with
$P^i = 0$ for $i > b$
(see for example the very general
Derived Categories, Lemma \ref{derived-lemma-subcategory-left-resolution}).
Note that $P^n$ is flat over $\mathcal{O}_{S, s}$ for all $n$.
Consider $\Coker(d_P^{a - 1})$. By More on Algebra, Lemma
\ref{more-algebra-lemma-last-one-flat}
this is a flat $\mathcal{O}_{S, s}$-module.
Hence by Lemma \ref{lemma-flat-up-down-henselization}
this is a flat $\mathcal{O}_{S, s}^{sh}$-module.
Thus $\tau_{\geq a}P^\bullet$ is a complex of
flat $\mathcal{O}_{S, s}^{sh}$-modules representing $M^{sh}$
in $D(\mathcal{O}_{S, s}^{sh}$ and we find that
$M^{sh}$ has tor amplitude in $[a, b]$, see
More on Algebra, Lemma \ref{more-algebra-lemma-tor-amplitude}.
Thus (1) implies (3). Of course this implies also
(2) $\Rightarrow$ (3) by replacing $\mathcal{O}_{S, s}$ by
$\mathcal{O}_{S, s}^h$. The same argument applies to prove
(4) $\Rightarrow$ (5).
\end{proof}

\end{document}
